\documentclass[11pt]{article}
\usepackage[margin=1in]{geometry}
\usepackage{amsmath,amssymb}
\usepackage{booktabs}
\usepackage{array}

\newcommand{\bR}{\mathbf{R}}
\newcommand{\bT}{\mathbf{T}}
\newcommand{\bs}{\mathbf{s}}
\newcommand{\bI}{\mathbf{I}}
\newcommand{\balpha}{\boldsymbol{\alpha}}
\newcommand{\bbeta}{\boldsymbol{\beta}}
\newcommand{\norm}[1]{\lVert #1 \rVert}

\title{Step clustering near the bottom of atmosphere\\
  in the adaptive Riccati solver}
\author{Technical note}
\date{April 2026}

\begin{document}
\maketitle

%% ================================================================
\section{Problem statement}

The forward Riccati solver \texttt{pydisort\_riccati\_jax} integrates
the coupled system
\begin{align}
  \frac{d\bR}{d\sigma}
    &= \balpha\bR + \bR\balpha + \bR\bbeta\bR + \bbeta,
    \label{eq:riccati}\\
  \frac{d\bT}{d\sigma}
    &= (\balpha + \bR\bbeta)\,\bT,
    \label{eq:T}\\
  \frac{d\bs}{d\sigma}
    &= (\balpha + \bR\bbeta)\,\bs + \bR\,\mathbf{q}_1 + \mathbf{q}_2
    \label{eq:s}
\end{align}
from $\sigma=0$ (BoA, $\tau=\tau_\mathrm{bot}$) to $\sigma=\tau_\mathrm{bot}$
(ToA, $\tau=0$), with initial conditions $\bR(0)=\mathbf{0}$,
$\bT(0)=\bI$, $\bs(0)=\mathbf{0}$.  The integration variable is
$\sigma$, mapped to optical depth as $\tau = \tau_\mathrm{bot}-\sigma$.
An adaptive Kvaerno5 solver (L-stable ESDIRK, order 5) with
PIDController step-size control determines the step grid.

\medskip
\noindent\textbf{Observation.}
The adaptive solver places roughly $22\text{--}28\%$ of all
steps in the region $\tau>25$ (near BoA), corresponding to
$\sigma<5$.  This clustering persists across a wide range of
tolerances ($10^{-2}$ to $5\times10^{-5}$).  This note
investigates why these steps arise and whether they are
necessary for the accuracy of the observable at ToA.


%% ================================================================
\section{Test problem}

We use the ``adiabatic cloud with drizzle'' profile:
$\tau_\mathrm{bot}=30$, $N_\mathrm{Quad}=16$ ($N=8$ half-streams),
$\mu_0=0.5$, $I_0=1$.  The single-scattering albedo and asymmetry
parameter are
\begin{align}
  \omega(\tau) &= 0.85 + 0.11\,\frac{\tau}{\tau_\mathrm{bot}}
    - 0.15\,\exp\!\Bigl[-\tfrac{1}{2}\Bigl(\frac{\tau-15}{0.5}\Bigr)^{\!2}\Bigr],\\
  g(\tau) &= 0.865 - 0.045\,\frac{\tau}{\tau_\mathrm{bot}}
    + 0.04\,\exp\!\Bigl[-\tfrac{1}{2}\Bigl(\frac{\tau-15}{0.5}\Bigr)^{\!2}\Bigr],
\end{align}
featuring a Gaussian perturbation (``drizzle spike'') centred at
$\tau=15$.  The Henyey--Greenstein phase function gives
Legendre coefficients $g_l(\tau) = g(\tau)^l$.


%% ================================================================
\section{Eigenvalue structure of $\balpha$}
\label{sec:eigenvalues}

The matrix $\balpha(\tau)$ entering the Riccati ODE is
\[
  \balpha = \mathbf{M}^{-1}\!\bigl(\omega\,\mathbf{D}^+\!\mathbf{W}
    - \bI\bigr),
\]
where $\mathbf{M}^{-1} = \mathrm{diag}(1/\mu_i)$ and
$\mathbf{D}^+\!\mathbf{W}$ is the discretised forward-scattering kernel.
Since $\omega\,\mathbf{D}^+\!\mathbf{W}$ has eigenvalues bounded by
$\omega<1$, the matrix $\balpha$ has all negative eigenvalues.

At the BoA ($\tau=30$, $\omega=0.96$) the eigenvalues of $\balpha$
for $N_\mathrm{Quad}=16$ are:

\begin{center}
\begin{tabular}{ccc}
\toprule
Stream $i$ & $\mu_i$ & $\lambda_i(\balpha)$\\
\midrule
1 & 0.020 & $-46.4$\\
2 & 0.102 & $-8.1$\\
3 & 0.237 & $-3.1$\\
4 & 0.408 & $-1.7$\\
5 & 0.592 & $-1.1$\\
6 & 0.763 & $-0.83$\\
7 & 0.898 & $-0.62$\\
8 & 0.980 & $-0.22$\\
\bottomrule
\end{tabular}
\end{center}

The stiffness ratio is $|\lambda_1/\lambda_8| \approx 214$,
almost entirely determined by $\mu_\mathrm{max}/\mu_\mathrm{min}
\approx 50$.  This stiffness is an intrinsic property of the
Gauss--Legendre quadrature: the steepest stream
($\mu_1\approx0.02$) sees an optical path $\tau/\mu_1\approx 50\tau$
per unit $\tau$, making it respond ${\sim}50$ times faster than
the most vertical stream ($\mu_8\approx0.98$).

The eigenvalues are nearly constant across $\tau$
(varying by $<1\%$) because the dominant $1/\mu_i$ factor in
$\mathbf{M}^{-1}$ overwhelms the $\tau$-dependence of
$\omega$ and $g$.


%% ================================================================
\section{Settling dynamics of $\bR$}

\subsection{Per-stream settling timescales}

Starting from $\bR(0)=\mathbf{0}$, the diagonal entries $R_{ii}$
settle to quasi-steady values on timescales proportional to $\mu_i$:

\begin{center}
\begin{tabular}{cccc}
\toprule
Stream $i$ & $\mu_i$ & $R_{ii}(\sigma=3)$ & $\sigma$ at 90\% of $R_{ii}(\sigma\!=\!3)$\\
\midrule
1 & 0.020 & $4.84\times10^{-2}$ & 0.045\\
2 & 0.102 & $8.98\times10^{-2}$ & 0.226\\
3 & 0.237 & $8.40\times10^{-2}$ & 0.799\\
4 & 0.408 & $7.29\times10^{-2}$ & 1.734\\
7 & 0.898 & $1.89\times10^{-2}$ & 2.729\\
8 & 0.980 & $7.05\times10^{-3}$ & 2.774\\
\bottomrule
\end{tabular}
\end{center}

The steepest stream ($\mu_1\approx0.02$) reaches 90\%
of its $\sigma=3$ value by $\sigma=0.045$ --- an interval much shorter
than a single step at the eventual ``cruising'' step size of
$\Delta\sigma\approx1$.  The shallowest stream
($\mu_8\approx0.98$) needs $\sigma\approx2.8$ to settle.

\subsection{Linear vs.\ nonlinear regime}

Decomposing the Riccati RHS into its linear part
$(\balpha\bR + \bR\balpha + \bbeta)$ and nonlinear part
$(\bR\bbeta\bR)$:

\begin{center}
\begin{tabular}{ccccl}
\toprule
$\sigma$ & $\norm{d\bR/d\sigma}$ & $\norm{\text{linear}}$ & $\norm{\bR\bbeta\bR}$ & Ratio\\
\midrule
$0.00$ & $1.05\times10^{1}$ & $1.05\times10^{1}$ & $0$ & $0$\\
$0.08$ & $1.00$ & $0.92$ & $0.13$ & $0.14$\\
$0.30$ & $0.30$ & $0.28$ & $0.23$ & $0.84$\\
$0.75$ & $0.13$ & $0.28$ & $0.30$ & $1.07$\\
$1.51$ & $0.070$ & $0.33$ & $0.34$ & $1.04$\\
$3.00$ & $0.031$ & $0.36$ & $0.36$ & $1.01$\\
\bottomrule
\end{tabular}
\end{center}

The nonlinear term $\bR\bbeta\bR$ becomes comparable to the
linear terms around $\sigma\approx0.5$.  After that, the two
contributions nearly cancel (the ratio $\approx 1$ while $\norm{d\bR/d\sigma}$
is small), indicating that $\bR$ has entered its quasi-steady
regime.

The transition from the linear-dominated to the cancellation regime
explains the step-size drop observed in the adaptive grid at
$\sigma\approx1.8$: the PIDController, having ramped up aggressively
during the linear phase, encounters the changed dynamics at the
nonlinear crossover and must recalibrate.


%% ================================================================
\section{Step distribution vs.\ tolerance}

Running the full solver at tolerances from $10^{-2}$ to $5\times10^{-5}$,
with the tightest tolerance as reference:

\begin{center}
\begin{tabular}{r r r r r r r}
\toprule
tol & Steps & BoA & Spike & ToA & $\max|\Delta u|$ & $\max|\Delta u/u|$\\
\midrule
$10^{-2}$          &  47 & 13 & 20 &  4 & $1.2\times10^{-3}$ & $3.8\times10^{-1}$\\
$5\times10^{-3}$   &  51 & 13 & 21 &  4 & $1.8\times10^{-3}$ & $1.4\times10^{-1}$\\
$10^{-3}$          &  69 & 15 & 31 &  4 & $9.7\times10^{-4}$ & $3.7\times10^{-2}$\\
$5\times10^{-4}$   &  86 & 19 & 39 &  5 & $2.9\times10^{-4}$ & $2.6\times10^{-2}$\\
$10^{-4}$          & 132 & 28 & 66 &  8 & $4.8\times10^{-5}$ & $1.6\times10^{-2}$\\
$5\times10^{-5}$   & 151 & 35 & 67 & 10 & \multicolumn{2}{c}{[reference]}\\
\bottomrule
\end{tabular}
\end{center}

The BoA fraction (steps with $\tau>25$) remains
$22\text{--}28\%$ at all tolerances.  The ToA radiance converges
as tolerance tightens, driven primarily by additional steps in the
spike region (20$\to$67) and near ToA (4$\to$10), with only
modest growth in BoA steps (13$\to$35).


%% ================================================================
\section{Contraction property: the BoA transient is forgotten}

The decisive test: inject a random symmetric perturbation
$\delta\bR$ with $\norm{\delta\bR}\approx 0.01$ to $\bR$
at various $\sigma$ values, then integrate to $\sigma=30$ (ToA)
and measure the perturbation in the final state.

\begin{center}
\begin{tabular}{c c c c}
\toprule
$\sigma_\mathrm{start}$ & $\norm{\delta\bR_\mathrm{init}}$
  & $\norm{\delta\bR_\mathrm{final}}$ & Attenuation factor\\
\midrule
$0.5$  & $6.1\times10^{-2}$ & $1.0\times10^{-8}$ & $1.7\times10^{-7}$\\
$1.0$  & $6.3\times10^{-2}$ & $4.4\times10^{-9}$ & $6.9\times10^{-8}$\\
$2.0$  & $4.4\times10^{-2}$ & $8.9\times10^{-9}$ & $2.0\times10^{-7}$\\
$5.0$  & $6.0\times10^{-2}$ & $1.8\times10^{-8}$ & $3.0\times10^{-7}$\\
$10.0$ & $5.1\times10^{-2}$ & $1.5\times10^{-7}$ & $2.9\times10^{-6}$\\
\bottomrule
\end{tabular}
\end{center}

Perturbations at $\sigma=0.5$ ($\tau=29.5$, near BoA) are
attenuated by a factor of ${\sim}10^7$ by the time they
reach ToA.  Even perturbations injected at $\sigma=10$
($\tau=20$, mid-atmosphere) are attenuated by ${\sim}10^6$.
The source vector $\bs$ shows the same insensitivity:
$\norm{\delta\bs_\mathrm{final}}\sim10^{-8}$.

This contraction is a consequence of the nonlinear Riccati
structure: the map $\bR_0 \mapsto \bR(\sigma)$ defined by
\eqref{eq:riccati} is a contraction for $\sigma$ large enough
that all modes have settled.  The settling timescale of the
slowest mode is ${\sim}1/|\lambda_8(\balpha)| \approx 5$,
so by $\sigma\approx 15$ (a few times this timescale), the
initial condition is forgotten to machine precision.


%% ================================================================
\section{Jacobian structure: direct sensitivity computation}
\label{sec:jacobian}

The contraction analysis (Section~5) provides indirect evidence
that the BoA region is informationally irrelevant.  We now
compute the Jacobians
$\partial u^+_0(\tau\!=\!0)/\partial\omega(\tau)$ and
$\partial u^+_0(\tau\!=\!0)/\partial g(\tau)$ directly via
JAX reverse-mode autodifferentiation (\texttt{jax.jacrev})
through the full Riccati ODE (forward and backward sweeps)
and BC solve.

\subsection{Method}

The atmospheric profiles $\omega(\tau)$ and
$g(\tau)$ (the Henyey--Greenstein asymmetry parameter,
with $g_l = g^l$) are perturbed by
piecewise-linear functions on a uniform grid of $K=61$ nodes
spanning $[0,\,\tau_\mathrm{bot}]$.  The Jacobians have shape
$(N,K) = (8,61)$: row~$i$ gives the sensitivity of
stream~$\mu_i$ at ToA to a perturbation at each $\tau$-node.

Reverse-mode AD is used because $N=8 < K=61$, requiring only
$N$ backward passes through the \texttt{RecursiveCheckpointAdjoint}
of \texttt{diffrax.diffeqsolve}.  The computation uses the
$m=0$ Fourier mode (azimuthally-averaged field, dominant
contribution); the Riccati contraction property applies
identically to all modes.

\subsection{Finite-difference validation}

Central finite differences ($\varepsilon = 10^{-5}$) at $\tau=0$
give $1{-}3\%$ agreement with the AD Jacobian.  At deeper $\tau$,
finite differences fail because the true sensitivity is so small
that the perturbation signal ($J\cdot\varepsilon \sim 10^{-8}$)
falls below the ODE solver noise floor --- itself a confirmation
of the extreme BoA insensitivity.

\subsection{Regional sensitivity}

We report $\norm{\partial u^+_0 / \partial p(\tau_j)}_2$
(column-wise $\ell_2$ norm across all $N=8$ streams) averaged
over each region:

\begin{center}
\begin{tabular}{l c c}
\toprule
Region & $\norm{\partial u^+_0/\partial\omega}_2$
       & $\norm{\partial u^+_0/\partial g}_2$\\
\midrule
ToA ($\tau<5$)       & $5.1\times10^{-2}$ & $5.3\times10^{-2}$\\
Mid ($5\!\leq\!\tau\!<\!13$) & $8.4\times10^{-4}$ & $6.2\times10^{-4}$\\
Spike ($13\!\leq\!\tau\!<\!17$) & $2.1\times10^{-5}$ & $1.6\times10^{-5}$\\
Deep ($17\!\leq\!\tau\!<\!25$) & $2.5\times10^{-6}$ & $1.0\times10^{-6}$\\
BoA ($\tau\geq 25$)  & $1.5\times10^{-7}$ & $7.4\times10^{-8}$\\
\bottomrule
\end{tabular}
\end{center}

The BoA-to-ToA mean sensitivity ratio is
$3.0\times10^{-6}$ for $\omega$ and $1.4\times10^{-6}$
for $g$, consistent with the ${\sim}10^{-6}$--$10^{-7}$
attenuation factors from the contraction test (Section~5).
Both parameters show the same spatial decay pattern:
sensitivity drops monotonically with depth, spanning five
orders of magnitude from ToA to BoA.

\subsection{Per-stream attenuation}

The attenuation is stream-dependent.  At $\tau=30$ (BoA)
versus $\tau=0$ (ToA), for the $\omega$ Jacobian:

\begin{center}
\begin{tabular}{c c c c c}
\toprule
Stream $i$ & $\mu_i$ & $|J_{i,\mathrm{BoA}}|$
  & $|J_{i,\mathrm{ToA}}|$ & Ratio\\
\midrule
1 & 0.020 & $5.3\times10^{-10}$ & $1.3\times10^{-1}$ & $4.2\times10^{-9}$\\
2 & 0.102 & $7.5\times10^{-10}$ & $1.3\times10^{-1}$ & $5.7\times10^{-9}$\\
3 & 0.237 & $1.1\times10^{-9}$ & $1.0\times10^{-1}$ & $1.0\times10^{-8}$\\
4 & 0.408 & $1.6\times10^{-9}$ & $6.8\times10^{-2}$ & $2.3\times10^{-8}$\\
5 & 0.592 & $2.3\times10^{-9}$ & $4.0\times10^{-2}$ & $5.8\times10^{-8}$\\
6 & 0.763 & $3.2\times10^{-9}$ & $2.6\times10^{-2}$ & $1.3\times10^{-7}$\\
7 & 0.898 & $4.3\times10^{-9}$ & $2.2\times10^{-2}$ & $1.9\times10^{-7}$\\
8 & 0.980 & $5.1\times10^{-9}$ & $1.4\times10^{-2}$ & $3.5\times10^{-7}$\\
\bottomrule
\end{tabular}
\end{center}

The steepest stream ($\mu_1\approx0.02$) is attenuated
by $4\times10^{-9}$, while the shallowest ($\mu_8\approx0.98$)
by $4\times10^{-7}$.  Steeper streams traverse more
optical path per unit $\tau$ (path $= \tau/\mu$), so
their signal is scattered more severely on the way from
BoA to ToA.

\subsection{Interpretation for the absorbing case}

A perturbation $\delta\omega = 0.01$ at BoA changes
$u^+_0(\mathrm{ToA})$ by ${\sim}10^{-9}$, which is
${\sim}2\times10^{-8}$ relative to $|u^+_0| \approx 0.05$.
This is far below any realistic measurement noise.
The ${\sim}15$ adaptive steps near BoA (22\% of the 69-step
grid) resolve a transient that is invisible to the ToA
observable.


%% ================================================================
\section{Conservative atmosphere: isolating the diffusion mechanism}
\label{sec:conservative}

The absorbing test problem (Sections~\ref{sec:jacobian}) combines
two mechanisms of information loss with depth: (i)~photon
absorption ($\omega < 1$), and (ii)~angular isotropisation
by multiple scattering (diffusion).  To disentangle them we
repeat the Jacobian computation for a \emph{conservative}
adiabatic cloud: $\omega = 1 - 10^{-6}$ (constant),
$g(\tau)$ linear from $g_\mathrm{top}=0.865$ to
$g_\mathrm{bot}=0.820$ (no drizzle spike), all other
parameters unchanged ($\tau_\mathrm{bot}=30$,
$N_\mathrm{Quad}=16$, $\mu_0=0.5$).  This simulates the
optical properties of an idealised cloud in the visible band.

\subsection{Grid distribution}

The adaptive solver selects 35~steps (vs.\ 69 in the absorbing
case): 5~at ToA ($\tau<5$), 14~at BoA ($\tau>25$, i.e.\ 40\%
of the grid).  The higher BoA fraction (40\% vs.\ 22\%) reflects
a slower settling transient in the near-conservative Riccati
equation, where the eigenvalues of~$\balpha$ are closer to zero.

\subsection{Regional sensitivity}

\begin{center}
\begin{tabular}{l c c}
\toprule
Region & $\norm{\partial u^+_0/\partial\omega}_2$
       & $\norm{\partial u^+_0/\partial g}_2$\\
\midrule
ToA ($\tau<5$)       & $5.7\times10^{-1}$ & $2.8\times10^{-2}$\\
Mid ($5\!\leq\!\tau\!<\!13$) & $3.3\times10^{-1}$ & $6.2\times10^{-3}$\\
Mid ($13\!\leq\!\tau\!<\!17$) & $2.0\times10^{-1}$ & $6.5\times10^{-3}$\\
Deep ($17\!\leq\!\tau\!<\!25$) & $9.4\times10^{-2}$ & $6.1\times10^{-3}$\\
BoA ($\tau\geq 25$)  & $2.0\times10^{-2}$ & $5.9\times10^{-3}$\\
\bottomrule
\end{tabular}
\end{center}

The results are qualitatively different from the absorbing case:

\begin{itemize}
\item \textbf{The $g$ Jacobian magnitude decays weakly.}
  The BoA-to-ToA sensitivity ratio is $2.1\times10^{-1}$
  (${\sim}5\times$ drop) --- far weaker than the
  five-order-of-magnitude decay in the absorbing atmosphere.
  However, the angular \emph{content} of the Jacobian collapses
  rapidly: by $\tau\approx 5$, all columns are proportional
  (Section~\ref{sec:angular}), so the relatively flat magnitude
  masks a loss of angular information.

\item \textbf{The $\omega$ Jacobian decays gently.}
  The BoA-to-ToA ratio is $3.6\times10^{-2}$
  (${\sim}28\times$ drop), five orders of magnitude weaker
  than the $3\times10^{-6}$ in the absorbing atmosphere.
  The residual decay arises because perturbing $\omega$
  introduces a small amount of absorption; the round-trip
  attenuation through $\tau=30$ of nearly conservative medium
  is modest but non-zero.
\end{itemize}

\subsection{Comparison: absorption vs.\ diffusion}

\begin{center}
\begin{tabular}{l l c c}
\toprule
Atmosphere & Parameter & BoA/ToA ratio & Decay mechanism\\
\midrule
Absorbing ($\omega\!\sim\!0.85{-}0.96$) & $\omega$
  & $3.0\times10^{-6}$ & absorption + diffusion\\
Absorbing & $g$
  & $1.4\times10^{-6}$ & absorption + diffusion\\[3pt]
Conservative ($\omega\!=\!1{-}10^{-6}$) & $\omega$
  & $3.6\times10^{-2}$ & residual absorption\\
Conservative & $g$
  & $2.1\times10^{-1}$ & diffusion (angular collapse)\\
\bottomrule
\end{tabular}
\end{center}

The five-order-of-magnitude decay in the absorbing case is
almost entirely due to photon absorption.  In the conservative
atmosphere, the $g$ Jacobian magnitude decays only modestly
($5\times$), but this masks a rapid loss of angular information:
by $\tau\approx 5$, all Jacobian columns are proportional
(rank~1), meaning the measurement cannot distinguish perturbations
at different depths (Section~\ref{sec:angular}).
In the language of diffusion theory, the diffusion length
diverges as $\omega\to 1$, so photon flux is preserved, but
the angular pattern freezes to a single diffusion eigenmode.


%% ================================================================
\section{Angular mode analysis: the HG-parameter Jacobian}
\label{sec:angular}

The $g$ Jacobian introduced in Section~\ref{sec:conservative}
decays only modestly in magnitude for conservative atmospheres.
To understand \emph{why} the BoA region is nevertheless
informationally redundant, we decompose the angular content
of $J_g$ using shifted Legendre polynomials.  The HG perturbation
$\delta g$ excites all Legendre orders via $\delta g_l = l\,g^{l-1}\,\delta g$
--- higher modes more strongly --- making it a broad probe of
angular diversity.

\subsection{Shifted-Legendre mode decomposition}

Each Jacobian column $J_g[:,\,j]$ (the 8-stream ToA response to a
$g$-perturbation at $\tau_j$) is projected onto shifted Legendre
polynomials $\tilde{P}_l(\mu) = P_l(2\mu - 1)$, orthogonal on
$[0,1]$ with the double-Gauss quadrature weights:
\[
  c_l(\tau_j) = (2l+1)\sum_{i=1}^{N}
    J_{g}[i,\,j]\;\tilde{P}_l(\mu_i)\;W_i.
\]

\begin{center}
\begin{tabular}{l c c c c c c}
\toprule
Region & $|c_0|$ & $|c_1|$ & $|c_2|$ & $|c_3|$ & $|c_4|$ & $|c_5|$\\
\midrule
ToA ($\tau<5$)  & $3.8\mathrm{e}{-3}$ & $7.3\mathrm{e}{-3}$
  & $7.5\mathrm{e}{-3}$ & $2.5\mathrm{e}{-3}$ & $4.6\mathrm{e}{-3}$ & $1.6\mathrm{e}{-2}$\\
Mid ($5\!\leq\!\tau\!<\!13$) & $2.1\mathrm{e}{-3}$ & $1.1\mathrm{e}{-3}$
  & $1.2\mathrm{e}{-4}$ & $6.6\mathrm{e}{-5}$ & $3.3\mathrm{e}{-5}$ & $1.7\mathrm{e}{-5}$\\
Deep ($17\!\leq\!\tau\!<\!25$) & $2.0\mathrm{e}{-3}$ & $1.1\mathrm{e}{-3}$
  & $1.2\mathrm{e}{-4}$ & $6.3\mathrm{e}{-5}$ & $3.3\mathrm{e}{-5}$ & $1.7\mathrm{e}{-5}$\\
BoA ($\tau\geq 25$)  & $2.0\mathrm{e}{-3}$ & $1.1\mathrm{e}{-3}$
  & $1.2\mathrm{e}{-4}$ & $6.1\mathrm{e}{-5}$ & $3.2\mathrm{e}{-5}$ & $1.7\mathrm{e}{-5}$\\
\bottomrule
\end{tabular}
\end{center}

Near ToA, multiple modes are active: $|c_2|\approx|c_1|$, and
$|c_5| > |c_1|$ (the HG perturbation excites higher modes more
strongly).  Below $\tau\approx 5$, the mode ratios freeze:
$|c_2/c_1| = 0.11$, $|c_3/c_1| = 0.058$, $|c_4/c_1| = 0.030$,
$|c_5/c_1| = 0.016$ --- constant from mid-atmosphere to BoA.
The angular pattern is not pure~$P_1$, but it is a \emph{single
fixed pattern} whose amplitude varies slowly with depth.

\subsection{Three definitions of retrieval-grid redundancy}

We evaluate three criteria for determining which $\tau$-grid
points are informationally redundant for the retrieval.

\paragraph{Definition A: BoA-anchored attribution.}
Starting from a fixed BoA point ($\tau=30$), greedily select the
column with the largest residual after projection onto the
current span.  This is column-pivoted QR with forced first pivot.

\begin{center}
\begin{tabular}{c c c c}
\toprule
Rank & $\tau$ selected & Residual & Region\\
\midrule
0 (anchor) & 30.0 & $3.0\times10^{-3}$ & BoA\\
1 & 0.0  & $1.7\times10^{-1}$ & ToA\\
2 & 1.0  & $1.6\times10^{-2}$ & ToA\\
3 & 0.5  & $4.9\times10^{-3}$ & ToA\\
4 & 1.5  & $1.2\times10^{-3}$ & ToA\\
5 & 2.0  & $6.4\times10^{-4}$ & ToA\\
6 & 3.5  & $1.0\times10^{-4}$ & ToA\\
7 & 2.5  & $2.3\times10^{-5}$ & ToA\\
\bottomrule
\end{tabular}
\end{center}

All $N=8$ selected points lie in $\tau\in[0,\,3.5]$.  No point
from $\tau>3.5$ is ever selected.  The cosine similarity of every
column at $\tau\geq 5$ with the BoA column is $\geq 0.9999$.

\paragraph{Definition B: Regional intrinsic dimensionality.}
SVD of $J_g$ restricted to depth regions:

\begin{center}
\begin{tabular}{l c c c c}
\toprule
Region & Cols & $\sigma_2/\sigma_1$ & Rank ($>$1\%) & Rank ($>$0.1\%)\\
\midrule
ToA ($\tau<5$)       & 10 & $1.4\times10^{-1}$ & 4 & 6\\
Mid ($5\!\leq\!\tau\!<\!25$)  & 40 & $2.3\times10^{-3}$ & 1 & 2\\
BoA ($\tau\geq 25$)  & 11 & $3.1\times10^{-8}$ & 1 & 1\\
\bottomrule
\end{tabular}
\end{center}

The BoA block is spectacularly rank-1 ($\sigma_2/\sigma_1 =
3\times10^{-8}$).  All 11 BoA columns are mutually proportional
to machine precision.  The mid block is also effectively rank-1.
Only the ToA block has multiple independent angular directions
(rank 4--6).

\paragraph{Definition C: Local resolving length.}
$L_\mathrm{res}(\tau) = \norm{J_g(\tau)} / \norm{dJ_g/d\tau}$
measures how finely the measurement can resolve $g(\tau)$:

\begin{center}
\begin{tabular}{l c c}
\toprule
Region & Mean $L_\mathrm{res}$ & Implied grid points\\
\midrule
ToA ($\tau<5$)       & 2.2  & 7.1\\
Mid ($5\!\leq\!\tau\!<\!13$)  & 3.3  & 15.8\\
Mid ($13\!\leq\!\tau\!<\!17$) & 1.0  & 7.9\\
Deep ($17\!\leq\!\tau\!<\!25$) & 2.5  & 4.5\\
BoA ($\tau\geq 25$)  & 8.9  & 3.3\\
\bottomrule
\end{tabular}
\end{center}

Definition~C prescribes ${\approx}16$ points in the mid
region where Definition~B gives rank~1.  This is the predicted
disagreement: the Jacobian amplitude varies locally
(short $L_\mathrm{res}$), but the angular \emph{pattern}
is frozen (rank~1).  Criterion~C detects amplitude variation
but misses global angular redundancy.

\subsection{Feature-reactivity test: Gaussian spike position sweep}

To test whether the angular collapse depends on the
$g(\tau)$ profile shape, we place a Gaussian spike
\[
  g_\mathrm{spike}(\tau) = g_\mathrm{smooth}(\tau)
  + 0.04\,\exp\!\Bigl[-\tfrac{1}{2}\Bigl(\frac{\tau-\tau_s}{0.5}\Bigr)^{\!2}\Bigr]
\]
at nine positions $\tau_s \in \{0,\,1,\,2,\,3,\,4,\,5,\,6,\,8,\,15\}$,
spanning from the surface boundary ($\tau=0$, half-spike) through the
collapse threshold ($\tau\approx 5$) and well past it.

\paragraph{Results.}
The angular structure is invariant across all spike positions:

\begin{center}
\begin{tabular}{c c c c c c}
\toprule
$\tau_s$ & Steps & ToA pts & ToA rank (1\%)
  & BoA $\sigma_2/\sigma_1$ & $\cos_\mathrm{sim}$ Mid min\\
\midrule
smooth  & 35 &  5 & 4 & $3.2\times10^{-8}$ & 0.9999\\
0 (half)& 50 & 20 & 4 & $3.4\times10^{-8}$ & 0.9999\\
1       & 57 & 27 & 4 & $3.6\times10^{-8}$ & 0.9999\\
2       & 60 & 30 & 4 & $3.6\times10^{-8}$ & 0.9999\\
3       & 54 & 24 & 4 & $3.6\times10^{-8}$ & 0.9999\\
4       & 56 & 23 & 4 & $3.6\times10^{-8}$ & 0.9999\\
5       & 57 & 15 & 4 & $3.6\times10^{-8}$ & 1.0000\\
6       & 57 &  5 & 4 & $3.6\times10^{-8}$ & 0.9999\\
8       & 55 &  4 & 4 & $3.6\times10^{-8}$ & 0.9999\\
15      & 54 &  4 & 4 & $3.7\times10^{-8}$ & 1.0000\\
\bottomrule
\end{tabular}
\end{center}

The frozen mode ratios are stable across all cases:
$|c_2/c_1| = 0.110{-}0.123$, $|c_3/c_1| = 0.058{-}0.069$.
The boundary half-spike ($\tau_s=0$) shifts the ratio
slightly ($|c_2/c_1|=0.123$) since the truncated Gaussian
modifies the effective surface optical properties, but the
rank structure is unchanged.
The greedy selection always picks
from $\tau<4$; no spike-region point is ever selected.

\paragraph{Solver response vs.\ angular invariance.}
Near-surface spikes produce the strongest solver response:
$\tau_s=2$ forces 30~ToA steps (vs.\ 5 for smooth), a
$6\times$ increase.  Yet the ToA rank, BoA rank, cosine
similarity profile, and rank structure are unchanged.
The spike alters the Jacobian \emph{amplitudes} near ToA
but not the angular \emph{pattern} --- the rank-4 subspace
at ToA is the same subspace regardless of spike position.

\paragraph{Interpretation.}
The angular collapse at $\tau\approx 5$ is a structural property
of multiple scattering in a thick medium, independent of the
profile shape.  Spikes at any depth --- including at the
surface boundary ($\tau_s=0$, half-spike) and before the
transition --- do not break or shift the collapse.
This profile-independence extends to the strongest possible
test: a spike \emph{at the measurement surface} ($\tau=0$).
The retrieval grid can therefore be prescribed once from the
angular structure, without knowing the specific $g(\tau)$ profile.

\subsection{Retrieval grid vs.\ solver grid}

The angular analysis reveals a fundamental mismatch between
the ODE solver's adaptive grid and the retrieval information
grid:

\begin{center}
\begin{tabular}{l r r}
\toprule
 & Smooth & Spike\\
\midrule
Solver steps (total)     & 35  & 54\\
Solver steps at BoA ($\tau>25$) & 14  & 14\\
Retrieval DoF at BoA (Def.\ B) & 1   & 1\\[3pt]
Solver steps near ToA ($\tau<5$) & 5   & 4\\
Retrieval DoF near ToA (Def.\ B) & 4--6 & 4--6\\
\bottomrule
\end{tabular}
\end{center}

Near BoA, the solver allocates 14~steps where the retrieval
needs 1 --- a $14\times$ over-resolution.  Near ToA, the solver
uses 4--5~steps where 4--6 independent quantities exist ---
marginal or under-resolved.  The solver grid is designed for ODE
accuracy (resolving the Riccati transient), not for retrieval
information content.  The retrieval parameterisation grid must
therefore be designed independently from the ODE step grid.


\subsection{Absorbing atmosphere: angular collapse with magnitude decay}

The preceding analysis used a conservative atmosphere to isolate
the diffusion mechanism.  We now compute $J_g$ for the absorbing
adiabatic cloud ($\omega \sim 0.85{-}0.96$ with drizzle spike at
$\tau=15$) to test whether the angular collapse co-locates with
the five-order-of-magnitude decay documented in
Section~\ref{sec:jacobian}.

\begin{center}
\begin{tabular}{l c c c}
\toprule
Region & $\norm{J_g}_2$ & $|c_2/c_1|$ & $\cos_\mathrm{sim}$ min\\
\midrule
ToA ($\tau<5$)       & $5.3\times10^{-2}$ & 0.621 & 0.36\\
Mid ($5\!\leq\!\tau\!<\!13$) & $6.2\times10^{-4}$ & 0.123 & 0.989\\
Spike ($13\!\leq\!\tau\!<\!17$) & $1.6\times10^{-5}$ & 0.286 & 1.000\\
Deep ($17\!\leq\!\tau\!<\!25$) & $1.0\times10^{-6}$ & 0.296 & 1.000\\
BoA ($\tau\geq 25$)  & $7.4\times10^{-8}$ & 0.297 & 1.000\\
\bottomrule
\end{tabular}
\end{center}

SVD of the BoA block gives $\sigma_2/\sigma_1 = 2.2\times10^{-5}$
(rank~1); the ToA block has rank~4 (1\%~threshold) and rank~6
(0.1\%), matching the conservative case.

Two differences from the conservative atmosphere are notable:
\begin{itemize}
\item The \emph{frozen mode ratios} differ: $|c_2/c_1| \approx 0.30$
  at depth (absorbing) vs.\ $0.11$ (conservative).  The angular
  pattern still collapses to a single fixed pattern, but the
  pattern itself depends on the absorption profile.
\item The \emph{angular collapse is slightly later}: $\cos_\mathrm{sim}$
  reaches 0.989 at the bottom of the mid region ($\tau\approx 13$)
  vs.\ $\geq 0.9999$ by $\tau\approx 5$ in the conservative case.
  The sliding-window rank in the mid region is 1.8 (absorbing)
  vs.\ 1.2 (conservative).
\end{itemize}
These differences do not affect the practical conclusion:
the ToA rank is identical (4--6), and below $\tau\approx 13$
the absorbing atmosphere is rank-1 with negligible absolute
sensitivity ($\norm{J_g} < 10^{-5}$).  In the absorbing case,
magnitude thresholding alone trims the grid more aggressively
than angular analysis.

\subsection{Full-radiance-field analysis}

The preceding analyses use only the zeroth Fourier mode
$u_0(\mu) = u_m(\mu)|_{m=0}$, which captures the azimuthally
averaged radiance.  The full radiance field
$u(\mu,\phi) = \sum_{m=0}^{N_\mathrm{Fourier}-1} u_m(\mu)\cos m(\phi_0-\phi)$
contains additional azimuthal information.  We test whether
this additional information increases the retrieval DoF by
stacking the per-mode Jacobians:
\[
  J_g^\mathrm{full} = \begin{pmatrix} J_g^{m=0} \\ J_g^{m=1} \\ \vdots \\ J_g^{m=M-1} \end{pmatrix}
  \in \mathbb{R}^{NM \times K}
\]
where $N = N_\mathrm{Quad}/2$ is the number of upwelling streams
and $M$ is the number of Fourier modes.

Three tiers of increasing measurement dimensionality were tested
at $\tau_s = 3$ (conservative, $\omega \approx 1$):

\begin{center}
\begin{tabular}{l c c c c c}
\toprule
Tier & $N_\mathrm{Quad}$ & Modes $M$ & Rows
  & ToA rank (1\%) & Mid rank (1\%)\\
\midrule
Baseline ($u_0$)  & 16 &  1 &   8 & 4 & 1\\
(a) Full $u$      & 16 & 16 & 128 & 4 & 2\\
(b) $u_0$, doubled& 32 &  1 &  16 & 4 & 1\\
\bottomrule
\end{tabular}
\end{center}

Neither azimuthal diversity (tier~a) nor doubled angular resolution
(tier~b) increases the ToA rank beyond~4.  Tier~(a) adds one degree
of freedom in the mid-atmosphere (rank $1\to 2$), but only where
$\norm{J_g} < 10^{-3}$ --- well below the ToA signal.
A third tier combining both ($N_\mathrm{Quad}=32$, all 32~Fourier
modes) was attempted but not completed due to computational cost
(each $\mathrm{jacrev}$ call at $N_\mathrm{Quad}=32$ exceeded
20~minutes per mode); the combined effect remains untested.

The per-mode BoA sensitivity reveals why azimuthal modes are
unlikely to help: higher Fourier modes decay exponentially with
depth.  At BoA ($\tau = 30$), $\norm{J_g^{m=1}} =
9.0\times10^{-12}$ and $\norm{J_g^{m=2}} = 1.2\times10^{-16}$, compared
to $\norm{J_g^{m=0}} = 3.1\times10^{-3}$.  Azimuthal modes are generated
by single scattering (which decays as $e^{-\tau/\mu_0}$) and destroyed
by multiple scattering, so they carry information only near the surface.
At depth, the field is isotropized and only $m=0$ survives.


\subsection{Single-scattering albedo scan}

Absorption delays the angular collapse
(Section~\ref{sec:angular}: $\tau\approx 13$ vs.\ $\tau\approx 5$)
but also reduces the Jacobian magnitude.  We scan
$\omega \in \{0.999,\, 0.99,\, 0.98,\, 0.95,\, 0.90\}$
(uniform with $\tau$, no spike) to determine whether a useful
window exists where the collapse is delayed but the signal remains
detectable.

\begin{center}
\begin{tabular}{c c c c c}
\toprule
$\omega$ & BoA $\norm{J_g}$ & ToA rank (1\%)
  & Mid sliding rank & $\cos_\mathrm{sim}$ Mid min\\
\midrule
$1 - 10^{-6}$ & $3.1\times10^{-3}$ & 4 & 1.2 & 1.0000\\
0.999          & $2.4\times10^{-3}$ & 4 & 1.1 & 1.0000\\
0.99           & $4.2\times10^{-4}$ & 4 & 1.5 & 0.9988\\
0.98           & $8.9\times10^{-5}$ & 4 & 1.6 & 0.9976\\
0.95           & $2.3\times10^{-6}$ & 4 & 1.6 & ---\\
0.90           & $2.2\times10^{-8}$ & 4 & 1.6 & ---\\
\bottomrule
\end{tabular}
\end{center}

The BoA $\norm{J_g}$ drops below $10^{-5}$ (the numerical noise
floor given ODE tolerance $10^{-3}$) between $\omega = 0.95$
and $\omega = 0.98$.  Absorption does delay the angular collapse:
the mid-atmosphere sliding rank increases from 1.2 (conservative)
to 1.6 ($\omega = 0.95$), and the cosine similarity departs from
unity.  However, the ToA rank remains~4 for all $\omega$ ---
absorption delays the collapse at depth but does not create new
independent directions where the signal is strong.

A full-u test at the threshold albedo ($\omega = 0.98$, spike at
$\tau_s = 3$) with all 16~Fourier modes at $N_\mathrm{Quad}=16$
gives ToA rank~4, matching the conservative result.  The
corresponding $N_\mathrm{Quad}=32$ test was not completed due to
computational cost.


\subsection{Interpretation}

\begin{enumerate}
\item The angular collapse is fast and complete in both regimes.
  In the conservative atmosphere, all higher angular modes
  decay to fixed fractions of the $P_1$ mode by
  $\tau\approx 5$ (roughly one transport mean free path
  $\ell^* = 1/(1-g) \approx 7$).
  In the absorbing atmosphere, the collapse occurs by
  $\tau\approx 13$, with the five-order magnitude decay
  providing a more aggressive trimming criterion.

\item Although $\norm{J_g}$ decays only modestly in the conservative
  atmosphere ($5\times$ from ToA to BoA), the angular structure
  collapses: all modes except $P_1$ are effectively zeroed by
  $\tau\approx 5$.  Magnitude alone is misleading --- the
  information content is determined by the angular rank.

\item The BoA region ($\tau\geq25$) is rank-1 to machine
  precision.  A perturbation $\delta g$ at any depth in this
  region produces the same angular signature at ToA --- the
  measurement cannot distinguish which depth the perturbation
  came from.  All 14 adaptive steps near BoA carry a single
  degree of freedom, attributable to one grid point.

\item Definitions A and B agree: the retrieval needs ${\sim}8$
  points in $\tau\in[0,3.5]$ and 1 point at BoA.  Definition~C
  disagrees in the mid-atmosphere, prescribing ${\sim}16$ points
  where only 1--2 independent quantities exist --- the
  ``individually relevant but collectively redundant'' phenomenon.

\item Spike tests at $\tau_\mathrm{spike}\in\{0,1,2,3,4,5,6,8,15\}$
  --- including a half-spike at the surface boundary ($\tau=0$) ---
  confirm that the angular collapse is insensitive to local
  profile features at \emph{every} depth.  The retrieval grid
  criteria (greedy selection, SVD ranks, mode ratios) are invariant
  across all spike positions.  Only the solver grid reacts to the
  feature (adding up to 25~steps for ODE accuracy).

\item When tested separately, neither azimuthal diversity
  (all 16~Fourier modes at $N_\mathrm{Quad}=16$) nor doubled
  angular resolution ($N_\mathrm{Quad}=32$, $m=0$ only) increases
  the ToA rank beyond~4.  Higher Fourier modes decay exponentially
  with depth ($\norm{J_g^{m=1}}$ at BoA is $10^{9}$ times smaller
  than $\norm{J_g^{m=0}}$), so they carry information only at the
  surface.  The combined test (all modes at $N_\mathrm{Quad}=32$)
  was not completed due to computational cost.

\item At $N_\mathrm{Quad}=16$, no useful $\omega$~window was found
  for exploiting delayed angular collapse.  Reducing $\omega$ from
  1 to 0.90 delays the collapse from $\tau\approx 5$ to past
  $\tau\approx 13$, but the BoA sensitivity simultaneously drops
  five orders of magnitude.  The ToA rank remains~4 for all
  $\omega$ tested.
\end{enumerate}


%% ================================================================
\section{Conclusions}

\begin{enumerate}
\item \textbf{The BoA step clustering is mathematically genuine.}
  It resolves a stiff initial transient of $\bR$ as the
  reflection matrix grows from zero.  The stiffness originates
  from the Gauss--Legendre quadrature: the steepest stream
  ($\mu_\mathrm{min}\approx 0.02$ for $N_\mathrm{Quad}=16$) has
  a settling timescale ${\sim}\mu_\mathrm{min}\approx 0.02$,
  which is ${\sim}50$ times shorter than the shallowest stream.
  The adaptive controller correctly resolves this fast
  transient.

\item \textbf{The T equation amplifies the clustering.}
  Integrating $\bR$ alone requires ${\sim}43$ steps vs.\
  ${\sim}50$ for the coupled $(\bR,\bT,\bs)$ system.  The
  fast eigenvalues of $\balpha$ drive exponential decay in
  $\bT$ (from $\bI$ down to $\norm{\bT}\approx 8\times10^{-4}$),
  adding ${\sim}7$ steps to track this decay.

\item \textbf{Retrieval information is concentrated near ToA
  in both absorption regimes
  (Sections~\ref{sec:jacobian}--\ref{sec:angular}).}
  \begin{itemize}
  \item \emph{Absorbing atmosphere}
    ($\omega \sim 0.85{-}0.96$): Jacobian magnitude drops
    five orders of magnitude from ToA to BoA
    (ratio~${\sim}10^{-6}$).  Angular collapse occurs by
    $\tau\approx 13$; magnitude thresholding alone suffices
    for grid trimming.
  \item \emph{Conservative atmosphere}
    ($\omega \approx 1$): Jacobian magnitude decays only
    modestly ($5\times$), but the angular content collapses
    by $\tau\approx 5$.  Below this depth, all Jacobian
    columns are rank-1 (proportional).  The BoA region is
    informationally redundant: 11 grid points carry a
    single degree of freedom.
  \end{itemize}
  Angular information collapses rapidly with depth in both
  regimes; the ToA region ($\tau < 5$) has rank~4--6, while
  deeper regions are rank-1.

\item \textbf{Implication for retrieval grid design.}
  The retrieval parameterisation grid for $g(\tau)$ should be:
  \begin{itemize}
  \item Dense in $\tau \in [0,\,{\sim}5]$: this is where all
    independent angular information resides (up to $N=8$
    degrees of freedom).
  \item Sparse below $\tau\approx 5$: only 1--2 grid points
    are needed (the information is one-dimensional).
  \item Always include a point at BoA ($\tau = \tau_\mathrm{bot}$)
    to anchor the single retrievable degree of freedom
    at depth.
  \end{itemize}
  This grid design applies to \emph{both} absorbing and
  conservative atmospheres.  The adaptive ODE grid is
  necessary for solver accuracy but is not the retrieval grid
  --- the retrieval grid should be designed from the Jacobian
  structure, not from the ODE step placement.

\item \textbf{The retrieval grid is profile-independent.}
  Gaussian spikes in $g(\tau)$ placed at nine positions
  ($\tau_s \in \{0,1,2,3,4,5,6,8,15\}$) --- including a
  half-spike at the surface boundary and spikes before, at,
  and after the angular collapse threshold --- leave all
  retrieval-grid criteria invariant.  The angular collapse at
  $\tau\approx 5$ is a structural property of multiple
  scattering in a thick medium.

\item \textbf{The full radiance field does not add retrieval DoF
  (tested separately).}
  The Jacobian was computed with all $N_\mathrm{Fourier}=16$
  azimuthal modes (at $N_\mathrm{Quad}=16$) and separately
  with doubled streams ($N_\mathrm{Quad}=32$, $m=0$ only).
  Neither increases the ToA rank beyond~4.  Higher Fourier
  modes ($m\geq 1$) decay exponentially with depth and carry
  information only at the surface.  The combined test (all
  modes at $N_\mathrm{Quad}=32$) was not completed due to
  computational cost.

\item \textbf{No useful single-scattering albedo window found
  (at $N_\mathrm{Quad}=16$).}
  Reducing $\omega$ from 1 to 0.90 delays the angular collapse
  from $\tau\approx 5$ to past $\tau\approx 13$, but the BoA
  sensitivity drops below the numerical noise floor ($10^{-5}$)
  between $\omega = 0.95$ and $\omega = 0.98$.  The ToA rank
  remains~4 for all~$\omega$ tested.  At the threshold
  ($\omega = 0.98$), a full-u test with 16~Fourier modes at
  $N_\mathrm{Quad}=16$ also gives rank~4; the $N_\mathrm{Quad}=32$
  test was not completed.
\end{enumerate}

\end{document}
